\chapter{अणुखंड कक्षक और बहुपरमाणुक अणु}\label{ch:b2:fragment-orbitals}

जल की लुइस संरचना दो तुल्य O--H आबंध और दो
तुल्य एकाकी युग्म दिखाती है। फिर भी जब जल पराबैंगनी प्रकाश से आयनित होता है, तो
इलेक्ट्रॉन चार भिन्न ऊर्जाओं के साथ निकलते हैं, दो नहीं; मेथेन, चार
तुल्य C--H आबंधों के साथ, दो दिखाती है। अणु में इलेक्ट्रॉन परमाणु-युग्मों के बीच
आबंधों में नहीं बैठते: वे पूरे अणु पर फैले कक्षकों में रहते हैं।
यह अध्याय उन कक्षकों को \omterm{def:b2:fragment-orbitals:fragment}{अणुखंडों} के कक्षकों से बनाता है, समझाता है कि
जल मुड़ा हुआ और एथीन समतल क्यों है, और कार्बधनायन अपने चारों ओर
ऐल्किल समूह क्यों पसंद करता है।

\begin{recall}
\Cref{ch:b2:diatomic-mos}: \omterm{def:b2:diatomic-mos:lcao}{LCAO}, अतिव्यापन और सममिति, दो अन्योन्यक्रियारत स्तर,
आबंधी और \omterm{def:b2:diatomic-mos:bonding}{प्रतिआबंधी कक्षक}। प्रथम वर्ष का खंड: VSEPR आकृतियाँ, वर्णन के रूप में
संकरण, अनुनाद, कार्बधनायन और उनके स्थायित्व का क्रम (रिक्त कक्षक के साथ
C--H आबंधों की अन्योन्यक्रिया की घोषणा वहीं की गई थी)।
\end{recall}

\section{अणुखंड विधि}

\begin{definition}[अणुखंड]\label{def:b2:fragment-orbitals:fragment}
\emph{अणुखंड}\index{अणुखंड} किसी अणु के परमाणुओं का वह समूह है जिस पर
अलग से विचार किया जाता है (एक परमाणु, हाइड्रोजनों का एक युग्म, एक \ce{CH2} समूह), और \emph{अणुखंड
कक्षक}\index{अणुखंड कक्षक} उसके कक्षकों में से एक। अणु के कक्षक
दो अणुखंडों के कक्षकों को, दो-दो करके, अन्योन्यक्रिया करने देकर बनाए जाते हैं,
जैसे द्विपरमाणुक अणु के परमाण्वीय कक्षक।
\end{definition}

\begin{definition}[सममिति-अनुकूलित संयोजन]\label{def:b2:fragment-orbitals:symmetry-adapted}
तुल्य परमाण्वीय कक्षकों का \emph{सममिति-अनुकूलित संयोजन}\index{सममिति-अनुकूलित संयोजन}
ऐसा संयोजन है जिसे अणु की हर सममिति संक्रिया (किसी दर्पण तल में
परावर्तन, किसी अक्ष के परितः घूर्णन) या तो अपरिवर्तित छोड़ती है या उसके विपरीत में बदल देती है,
या, अपभ्रष्ट समुच्चयों के लिए, उसी समुच्चय के संयोजनों में
रूपांतरित करती है।
\end{definition}

\begin{proposition}[दो हाइड्रोजन]\label{prop:b2:fragment-orbitals:h2-pair}
1s कक्षकों $h_1$ और $h_2$ वाले दो तुल्य हाइड्रोजनों के लिए
\omterm{def:b2:fragment-orbitals:symmetry-adapted}{सममिति-अनुकूलित संयोजन} $h_1 + h_2$ और $h_1 - h_2$ हैं: दोनों परमाणुओं का आदान-प्रदान करने वाली
किसी भी संक्रिया में पहला अपरिवर्तित रहता है, दूसरा चिह्न बदलता है।
\end{proposition}

\begin{proof}
परमाणुओं का आदान-प्रदान करने वाली संक्रिया $h_1$ को $h_2$ में और $h_2$ को
$h_1$ में बदलती है: $h_1 + h_2$ अपरिवर्तित रहता है और $h_1 - h_2$, $h_2 - h_1 = -(h_1 - h_2)$ बन जाता है।
\end{proof}

\begin{proposition}[तीन और चार हाइड्रोजन]\label{prop:b2:fragment-orbitals:h4}
किसी चतुष्फलक के कोनों पर चार हाइड्रोजनों के लिए संयोजन हैं: एक
पूर्णतः सममित संयोजन $h_1 + h_2 + h_3 + h_4$ और तीन अपभ्रष्ट
संयोजनों का एक समुच्चय, हर एक में एक नोडीय तल, p कक्षकों जैसी आकृति वाले। किसी त्रिभुज के तीन
हाइड्रोजनों के लिए (जैसे अमोनिया में), एक सममित संयोजन और एक
अपभ्रष्ट युग्म।
\end{proposition}

\admitted

अपभ्रष्ट संयोजनों की आकृतियाँ, और उनके नाम (a, e, t), समूह सिद्धांत से आते हैं,
जो तृतीय वर्ष के खंड में विकसित किया गया है; यहाँ नाम केवल चिह्नक हैं।

\begin{proposition}[केवल समान, समान से अन्योन्यक्रिया करता है]\label{prop:b2:fragment-orbitals:interaction-rules}
कोई \omterm{def:b2:fragment-orbitals:fragment}{अणुखंड कक्षक} दूसरे \omterm{def:b2:fragment-orbitals:fragment}{अणुखंड} के केवल उन्हीं कक्षकों से अन्योन्यक्रिया करता है जो
अणु की हर सममिति संक्रिया में उसी तरह व्यवहार करते हैं; अन्योन्यक्रिया की
प्रबलता अतिव्यापन के साथ बढ़ती है और ऊर्जा अंतराल के साथ
घटती है।
\end{proposition}

\begin{proof}
यदि दो कक्षक किसी परावर्तन में भिन्न व्यवहार करते हैं, तो उस तल के सापेक्ष एक सममित और
दूसरा प्रतिसममित है, और उनका अतिव्यापन तथा
\omterm{def:b2:diatomic-mos:integrals}{अनुनाद समाकल} शून्य होते हैं (\cref{prop:b2:diatomic-mos:symmetry})। अतिव्यापन और
अंतराल पर निर्भरता \cref{thm:b2:diatomic-mos:heteronuclear} वाली है।
\end{proof}

\begin{method}[अणुखंड आरेख]\label{met:b2:fragment-orbitals:fragment-diagram}
\begin{enumerate}
  \item अणु को दो \omterm{def:b2:fragment-orbitals:fragment}{अणुखंडों} में बाँटिए, प्रायः एक केंद्रीय परमाणु और उसके चारों ओर के
        परमाणुओं का समुच्चय।
  \item बाहरी कक्षकों के \omterm{def:b2:fragment-orbitals:symmetry-adapted}{सममिति-अनुकूलित संयोजन} बनाइए।
  \item दोनों \omterm{def:b2:fragment-orbitals:fragment}{अणुखंडों} के कक्षकों को अणु के दर्पण तलों और अक्षों के
        सापेक्ष उनके व्यवहार के अनुसार छाँटिए।
  \item हर कक्षक को समान व्यवहार वाले कक्षकों से अन्योन्यक्रिया करने दीजिए, ऊर्जा में
        निकट होने पर अधिक प्रबलता से; बचे हुए अनाबंधी रहते हैं।
  \item संयोजकता इलेक्ट्रॉनों से भरिए और पढ़िए: उच्चतम भरा स्तर,
        आबंधी लक्षण, आकृति की वरीयताएँ।
\end{enumerate}
\end{method}

\section{जल}

जल को $yz$ तल में रखिए, $z$ अक्ष H--O--H कोण को समद्विभाजित करता हुआ। दो
दर्पण तल $z$ अक्ष को समाविष्ट करते हैं: आण्विक तल $yz$ और उसके लंबवत
तल $xz$। $xz$ तल के सापेक्ष, जो दोनों हाइड्रोजनों का आदान-प्रदान
करता है, $h_1 + h_2$ सममित है और $h_1 - h_2$ प्रतिसममित। ऑक्सीजन पर 2s और
$2\mathrm p_z$ दोनों तलों के सापेक्ष सममित हैं; $2\mathrm p_y$ (आण्विक
तल में, $z$ के लंबवत) $xz$ के सापेक्ष प्रतिसममित है;
$2\mathrm p_x$ आण्विक तल के सापेक्ष प्रतिसममित है, जो कोई भी
हाइड्रोजन संयोजन नहीं है।

अतः $h_1 + h_2$, 2s और $2\mathrm p_z$ के साथ मिश्रित होता है, $h_1 - h_2$, $2\mathrm p_y$ के साथ, और
$2\mathrm p_x$ अनाबंधी रहता है। छह कक्षक बनते हैं, जिनमें से चार आठ
संयोजकता इलेक्ट्रॉनों से भरे होते हैं।

\begin{omfigure}
\begin{tikzpicture}[font=\small]
  % O fragment (left)
  \pic (os) at (0,-2.4) {omlevel={ud}{0.7}}; \node[left] at (-0.45,-2.4) {2s};
  \pic (op1) at (-0.45,0) {omlevel={ud}{0.4}}; \pic (op2) at (0,0) {omlevel={u}{0.4}}; \pic (op3) at (0.45,0) {omlevel={u}{0.4}};
  \node[left] at (-0.7,0) {2p};
  \node at (0,-3.2) {O};
  % H2 fragment (right)
  \pic (hp) at (6,0.9) {omlevel={u}{0.7}}; \node[right] at (6.4,0.9) {$h_1 + h_2$};
  \pic (hm) at (6,1.8) {omlevel={u}{0.7}}; \node[right] at (6.4,1.8) {$h_1 - h_2$};
  \node[anchor=north west] at (6.4,0.6) {\ce{H2} अणुखंड};
  % MOs
  \pic (m1) at (3,-2.9) {omlevel={ud}{0.7}};  \node[omsymlabel, below=0.17cm, inner sep=0.5pt] at (m1-c) {$2a_1$};
  \pic (m2) at (3,-1.55) {omlevel={ud}{0.7}};  \node[omsymlabel, below=0.17cm, inner sep=0.5pt] at (m2-c) {$1b_2$};
  \pic (m3) at (3,-0.78) {omlevel={ud}{0.7}};  \node[omsymlabel, below=0.17cm, inner sep=0.5pt] at (m3-c) {$3a_1$};
  \pic (m4) at (3,0) {omlevel={ud}{0.7}};  \node[omsymlabel, below=0.17cm, inner sep=0.5pt] at (m4-c) {$1b_1$};
  \pic (m5) at (3,2.4) {omlevel={empty}{0.7}}; \node[omsymlabel, below=0.17cm, inner sep=0.5pt] at (m5-c) {$4a_1$};
  \pic (m6) at (3,3.1) {omlevel={empty}{0.7}}; \node[omsymlabel, below=0.17cm, inner sep=0.5pt] at (m6-c) {$2b_2$};
  \draw[omcorr, densely dashed] (os-r) -- (m1-l) (op3-r) -- (m2-l) (op3-r) -- (m3-l) (op3-r) -- (m4-l) (op3-r) -- (m5-l) (op3-r) -- (m6-l);
  \draw[omcorr, densely dashed] (hp-l) -- (m1-r) (hm-l) -- (m2-r) (hp-l) -- (m3-r) (hp-l) -- (m5-r) (hm-l) -- (m6-r);
\end{tikzpicture}

{\small जल का \omterm{def:b2:fragment-orbitals:fragment}{अणुखंड} आरेख, व्यवस्था-रूप। हाइड्रोजनों का सममित संयोजन
ऑक्सीजन के 2s और $2\mathrm p_z$ कक्षकों ($a_1$ कक्षक) के साथ मिश्रित होता है,
प्रतिसममित संयोजन $2\mathrm p_y$ ($b_2$) के साथ; आण्विक तल के लंबवत $2\mathrm p_x$
अनाबंधी रहता है ($1b_1$, ऑक्सीजन के 2p के स्तर पर)। चार भरे संयोजकता
स्तर, चार प्रकाशइलेक्ट्रॉन बैंड; सबसे ऊँचा, $1b_1$,
\qty{12.62}{eV} पर आयनित होता है। चिह्नक
$a_1$, $b_1$, $b_2$ नाम हैं (समूह सिद्धांत से, तृतीय वर्ष के खंड में)।}
% ledger: ie:H2O
\end{omfigure}

\begin{proposition}[जल मुड़ा हुआ क्यों है]\label{prop:b2:fragment-orbitals:bent-water}
रैखिक H--O--H में अक्ष के लंबवत ऑक्सीजन के दो p कक्षक
अनाबंधी होते हैं। मोड़ने से उनमें से एक, जो आण्विक तल में है,
$h_1 + h_2$ से अतिव्यापित होकर 2s कक्षक के साथ मिश्रित हो पाता है: वह भरा स्तर नीचे गिरता है, उससे अधिक जितना
अन्य भरे स्तर ऊपर उठते हैं। आठ संयोजकता इलेक्ट्रॉनों के साथ मुड़ा हुआ अणु
अधिक स्थायी है।
\end{proposition}

\begin{proof}[तर्क]
रैखिक अणु में ($z$, H--O--H के अनुदिश) $h_1 + h_2$ का
$2\mathrm p_x$ और $2\mathrm p_y$ से अतिव्यापन शून्य है (हर एक अक्ष को समाविष्ट करने वाले किसी तल के सापेक्ष
प्रतिसममित है, संयोजन सममित); वे दो अपभ्रष्ट
अनाबंधी युग्म हैं। $yz$ तल में मोड़ने पर $2\mathrm p_y$ आंशिक रूप से
हाइड्रोजनों की ओर इंगित करता है: $h_1 + h_2$ के साथ उसका अतिव्यापन शून्य से बढ़ता है, और वह भरा स्तर
जिसे वह वहन करता है स्थायी होता है (मुड़े अणु का $3a_1$ स्तर)। आबंधी
$1b_2$ स्तर थोड़ा अतिव्यापन खोता है और थोड़ा ऊपर उठता है। शुद्ध परिणाम कमी है,
जो तब सबसे बड़ी होती है जब नीचे जाता स्तर भरा हो, जैसा आठ इलेक्ट्रॉनों के साथ है।
केवल चार संयोजकता इलेक्ट्रॉनों वाले अणु, जैसे \ce{BeH2}, में वह स्तर
खाली होता है और वह रैखिक है।
\end{proof}

\section{मेथेन, अमोनिया और प्रकाशइलेक्ट्रॉन स्पेक्ट्रम}

मेथेन में चार हाइड्रोजन कक्षक सममित संयोजन देते हैं, जो
कार्बन के 2s कक्षक से मेल खाता है, और तीन का एक अपभ्रष्ट समुच्चय, जो
उसके तीन 2p कक्षकों से मेल खाता है। चार \omterm{def:b2:diatomic-mos:bonding}{आबंधी कक्षक} भरे होते हैं: एक नीचा ($a_1$) और तीन
उच्चतर ऊर्जा पर अपभ्रष्ट ($t_2$)। अमोनिया, तीन हाइड्रोजनों के साथ, नाइट्रोजन पर
हाइड्रोजनों से दूर इंगित करता एक भरा कक्षक, अपना एकाकी युग्म, उच्चतम भरे स्तर के रूप में
रखती है।

\begin{omfigure}
\begin{tikzpicture}[font=\small]
  \pic (cs) at (0,-1.6) {omlevel={ud}{0.7}}; \node[left] at (-0.45,-1.6) {2s};
  \pic (cp1) at (-0.45,0) {omlevel={u}{0.4}}; \pic (cp2) at (0,0) {omlevel={u}{0.4}}; \pic (cp3) at (0.45,0) {omlevel={empty}{0.4}};
  \node[left] at (-0.7,0) {2p};
  \node at (0,-2.5) {C};
  \pic (ha) at (6,-0.8) {omlevel={ud}{0.7}}; \node[right] at (6.4,-0.8) {$a_1$};
  \pic (ht1) at (5.55,0.4) {omlevel={u}{0.4}}; \pic (ht2) at (6,0.4) {omlevel={u}{0.4}}; \pic (ht3) at (6.45,0.4) {omlevel={empty}{0.4}};
  \node[right] at (6.7,0.4) {$t_2$};
  \node at (6,-1.6) {\ce{H4} अणुखंड};
  \pic (m1) at (3,-2.4) {omlevel={ud}{0.7}}; \node[omsymlabel, below=0.17cm, inner sep=0.5pt] at (m1-c) {$1a_1$};
  \pic (m2a) at (2.55,-0.9) {omlevel={ud}{0.4}}; \pic (m2b) at (3,-0.9) {omlevel={ud}{0.4}}; \pic (m2c) at (3.45,-0.9) {omlevel={ud}{0.4}};
  \node[omsymlabel, below=0.17cm, inner sep=0.5pt] at (m2b-c) {$1t_2$};
  \pic (m3) at (3,1.8) {omlevel={empty}{0.7}}; \node[omsymlabel, below=0.17cm, inner sep=0.5pt] at (m3-c) {$2a_1$};
  \pic (m4a) at (2.55,2.5) {omlevel={empty}{0.4}}; \pic (m4b) at (3,2.5) {omlevel={empty}{0.4}}; \pic (m4c) at (3.45,2.5) {omlevel={empty}{0.4}};
  \node[omsymlabel, below=0.17cm, inner sep=0.5pt] at (m4b-c) {$2t_2$};
  \draw[omcorr, densely dashed] (cs-r) -- (m1-l) (cp3-r) -- (m2a-l) (cs-r) -- (m3-l) (cp3-r) -- (m4a-l);
  \draw[omcorr, densely dashed] (ha-l) -- (m1-r) (ht1-l) -- (m2c-r) (ha-l) -- (m3-r) (ht1-l) -- (m4c-r);
\end{tikzpicture}

{\small मेथेन का \omterm{def:b2:fragment-orbitals:fragment}{अणुखंड} आरेख, व्यवस्था-रूप। चार हाइड्रोजनों का सममित संयोजन
कार्बन के 2s कक्षक से मेल खाता है, तीन अपभ्रष्ट संयोजन उसके 2p
कक्षकों से: दो भरे स्तर, अतः दो प्रकाशइलेक्ट्रॉन बैंड, जिनमें ऊपर वाला
(पहले आयनित, \qty{12.61}{eV} पर) छह इलेक्ट्रॉन रखता है।}
% ledger: ie:CH4
\end{omfigure}

\begin{definition}[प्रकाशइलेक्ट्रॉन स्पेक्ट्रमिकी]\label{def:b2:fragment-orbitals:photoelectron}
\emph{प्रकाशइलेक्ट्रॉन स्पेक्ट्रमिकी}\index{प्रकाशइलेक्ट्रॉन स्पेक्ट्रमिकी} ज्ञात ऊर्जा के
एकवर्णी विकिरण द्वारा अणुओं से निकाले गए इलेक्ट्रॉनों की गतिज ऊर्जाएँ
मापती है; स्पेक्ट्रम का हर बैंड किसी एक भरे स्तर से इलेक्ट्रॉन हटाने के
संगत है, और उसकी स्थिति उस स्तर की आयनन
ऊर्जा देती है।
\end{definition}

\begin{proposition}[कूपमान्स सन्निकटन]\label{prop:b2:fragment-orbitals:koopmans}
किसी भरे कक्षक से इलेक्ट्रॉन हटाने के लिए आवश्यक ऊर्जा
लगभग उस कक्षक की ऊर्जा का ऋण है, $I \approx -\varepsilon$।
\end{proposition}

\admitted

यह सन्निकटन आयनन के बाद अन्य इलेक्ट्रॉनों के शिथिलन की
उपेक्षा करता है; इसका औचित्य तृतीय वर्ष के खंड में दिया गया है। यह प्रकाशइलेक्ट्रॉन
स्पेक्ट्रम को कक्षक आरेख में बदल देता है: जल चार संयोजकता बैंड दिखाता है, मेथेन दो,
\omterm{def:b2:fragment-orbitals:fragment}{अणुखंड} आरेखों के अनुरूप --- और जल में दो तुल्य
एकाकी युग्मों और दो तुल्य आबंधों के चित्र के विपरीत, जो उसी कुल
इलेक्ट्रॉन घनत्व का दूसरे ढंग से वर्णन करता है परंतु हटाए गए इलेक्ट्रॉनों की
ऊर्जाएँ नहीं देता।

\section{एथीन}

एथीन को दो \ce{CH2} \omterm{def:b2:fragment-orbitals:fragment}{अणुखंडों} से बनाया जा सकता है। हर \ce{CH2} के पास, अपने दो C--H आबंधों के
कक्षकों के अतिरिक्त, तल में हाइड्रोजनों से दूर इंगित करता एक कक्षक
(जो अपने साथी के साथ C--C $\sigma$ आबंध बनाता है) और तल के लंबवत एक p कक्षक
होता है। दोनों p कक्षक संयोजित होकर एक भरा $\pi$ और एक
खाली $\pi^*$ बनाते हैं।

\begin{omfigure}
\begin{tikzpicture}[font=\small]
  \begin{scope}
    \pic at (0,0) {omp={0.85}{90}}; \pic at (1.2,0) {omp={0.85}{90}};
    \draw[thick] (0,0) -- (1.2,0);
    \draw[thick] (0,0) -- (-0.5,-0.45) (0,0) -- (-0.55,0.35) (1.2,0) -- (1.7,-0.45) (1.2,0) -- (1.75,0.35);
    \node at (0.6,-1.1) {$\pi$ (भरा)};
  \end{scope}
  \begin{scope}[xshift=3.6cm]
    \pic at (0,0) {omp={0.85}{90}}; \pic at (1.2,0) {omp={-0.85}{90}};
    \draw[thick] (0,0) -- (1.2,0);
    \draw[thick] (0,0) -- (-0.5,-0.45) (0,0) -- (-0.55,0.35) (1.2,0) -- (1.7,-0.45) (1.2,0) -- (1.75,0.35);
    \node at (0.6,-1.1) {$\pi^*$ (खाली)};
  \end{scope}
  \begin{scope}[xshift=7.4cm]
    \pic at (0,0) {omp={0.85}{90}}; \pic at (1.2,0) {omp={0.85}{0}};
    \draw[thick] (0,0) -- (1.2,0);
    \node at (0.6,-1.1) {$90^\circ$ से ऐंठा हुआ};
    \node[font=\scriptsize, align=center] at (0.6,1.3) {शून्य अतिव्यापन};
  \end{scope}
\end{tikzpicture}

{\small एथीन का $\pi$ निकाय, व्यवस्था-रूप (अणु किनारे से देखा गया,
हाइड्रोजन तल में)। तल के लंबवत दोनों p कक्षक समकला में
($\pi$, भरा) और विषमकला में ($\pi^*$, खाली) संयोजित होते हैं। एक \ce{CH2} को
$90^\circ$ से ऐंठने पर दोनों p कक्षक लंबवत हो जाते हैं: उनका अतिव्यापन, और $\pi$
आबंध, लुप्त हो जाते हैं।}
\end{omfigure}

$\pi$ आबंध एथीन को समतल रखता है: एक सिरे को ऐंठने पर p कक्षक एक-दूसरे से
दूर मुड़ जाते हैं, उनका अतिव्यापन ऐंठन कोण के कोज्या के अनुसार घटता है, और $90^\circ$ पर
$\pi$ आबंध समाप्त हो जाता है। अणु को ऐंठने के लिए $\pi$ आबंध की ऊर्जा देनी पड़ती है,
इसीलिए ऐल्कीनों के सिस और ट्रांस समावयवी कमरे के ताप पर एक-दूसरे में नहीं बदलते,
जबकि एकल आबंध के परितः घूर्णन मुक्त है।

\section{अतिसंयुग्मन}

\begin{definition}[अतिसंयुग्मन]\label{def:b2:fragment-orbitals:hyperconjugation}
\emph{अतिसंयुग्मन}\index{अतिसंयुग्मन} किसी C--H या C--C आबंध के भरे
$\sigma$ कक्षक की उसके समांतर किसी पड़ोसी खाली या आंशिक रूप से भरे p
या $\pi^*$ कक्षक के साथ अन्योन्यक्रिया है।
\end{definition}

\begin{proposition}[ऐल्किल समूह कार्बधनायनों को स्थायी करते हैं]\label{prop:b2:fragment-orbitals:hyperconjugation}
कार्बधनायन पड़ोसी कार्बन के हर उस C--H या C--C आबंध से स्थायी होता है जो
उसके खाली p कक्षक के साथ संरेखित हो सके: ऐसे आबंध जितने अधिक, धनायन उतना
अधिक स्थायी, जिससे क्रम तृतीयक > द्वितीयक > प्राथमिक > मेथिल निकलता है।
\end{proposition}

\begin{proof}[तर्क]
धनायनी कार्बन का खाली p कक्षक और उसके पास का भरा $\sigma_{\text{C--H}}$ कक्षक
बड़े अंतराल $\Delta$ और छोटे युग्मन
$\beta$ (शून्येतर केवल तब जब आबंध p कक्षक के लगभग समांतर हो) वाला द्वि-स्तरीय निकाय बनाते हैं।
\cref{thm:b2:diatomic-mos:heteronuclear} से भरा $\sigma$ स्तर लगभग
$\beta^2/\Delta$ नीचे आता है, और दोनों इलेक्ट्रॉन उसका दुगुना पाते हैं; खाली स्तर ऊपर उठता है,
बिना किसी लागत के। हर उपयुक्त रूप से संरेखित आबंध अपना स्थायीकरण जोड़ता है; मेथिल समूह
का एक C--H आबंध, उसके घूर्णन चाहे जो हो, सदा लगभग संरेखित रहता है।
\end{proof}

\begin{omfigure}
\begin{tikzpicture}[font=\small]
  \pic at (0,0) {omp={0.9}{90}};
  \node[circle, fill=black, inner sep=1.2pt] at (0,0) {};
  \node[font=\scriptsize, anchor=north west] at (0.08,-0.05) {C$^+$};
  \draw[thick] (0,0) -- (1.4,0);
  \node[circle, fill=black, inner sep=1.2pt] at (1.4,0) {};
  \draw[thick] (1.4,0) -- (1.75,0.85);
  \pic at (1.58,0.42) {omlobe={0.75}{68}};
  \pic at (1.58,0.42) {omlobe={0.45}{248}};
  \node[font=\scriptsize] at (1.95,1.0) {H};
  \draw[thick] (1.4,0) -- (2.0,-0.45) (1.4,0) -- (1.1,-0.6);
  \draw[<->, omProp, thick] (0.35,0.8) .. controls (0.8,1.2) and (1.2,1.2) .. (1.5,0.95);
  \node[font=\scriptsize, align=left, anchor=west] at (2.4,0.4) {खाली p कक्षक के समांतर\\ भरा $\sigma_{\text{C--H}}$ पालि};
  % level scheme
  \begin{scope}[xshift=7.2cm, yshift=-0.4cm]
    \pic (p) at (0,1.4) {omlevel={empty}{0.7}}; \node[left] at (-0.4,1.4) {खाली p};
    \pic (s) at (2.2,-0.8) {omlevel={ud}{0.7}}; \node[right] at (2.6,-0.8) {$\sigma_{\text{C--H}}$};
    \pic (lo) at (1.1,-1.2) {omlevel={ud}{0.7}};
    \pic (hi) at (1.1,1.7) {omlevel={empty}{0.7}};
    \draw[omcorr, densely dashed] (p-r) -- (lo-l) (p-r) -- (hi-l) (s-l) -- (lo-r) (s-l) -- (hi-r);
    \node[font=\scriptsize, anchor=north] at (1.1,-1.35) {$\approx \beta^2/\Delta$ से स्थायीकृत};
  \end{scope}
\end{tikzpicture}

{\small कार्बधनायन में \omterm{def:b2:fragment-orbitals:hyperconjugation}{अतिसंयुग्मन}, व्यवस्था-रूप। पड़ोसी कार्बन पर
भरा C--H \omterm{def:b2:diatomic-mos:bonding}{आबंधी कक्षक}, खाली p कक्षक के समांतर, अपने
इलेक्ट्रॉन घनत्व का कुछ भाग उसे दे देता है; भरा स्तर भिन्न ऊर्जाओं वाले दो स्तरों की
किसी भी अन्योन्यक्रिया की तरह स्थायी होता है।}
\end{omfigure}

\section{अभ्यास}

\begin{exercise}[$\star$]\label{exo:b2:fragment-orbitals:1}
जल के हाइड्रोजनों के 1s कक्षकों के दोनों \omterm{def:b2:fragment-orbitals:symmetry-adapted}{सममिति-अनुकूलित संयोजन} लिखिए
और बताइए कि उनमें से कौन उस तल में चिह्न बदलता है जो H--O--H कोण को
अणु के लंबवत समद्विभाजित करता है।
\end{exercise}

\begin{exercise}[$\star$]\label{exo:b2:fragment-orbitals:2}
जल के कितने संयोजकता \omterm{def:b2:diatomic-mos:lcao}{आण्विक कक्षक} हैं, कितने संयोजकता इलेक्ट्रॉन,
और कितने भरे संयोजकता कक्षक?
\end{exercise}

\begin{exercise}[$\star$]\label{exo:b2:fragment-orbitals:3}
मेथेन का प्रकाशइलेक्ट्रॉन स्पेक्ट्रम दो संयोजकता बैंड क्यों दिखाता है, जिनमें से एक
दूसरे से तीन गुना तीव्र है?
\end{exercise}

\begin{exercise}[$\star$]\label{exo:b2:fragment-orbitals:4}
एथीन के छहों परमाणु एक तल में क्यों हैं?
\end{exercise}

\begin{exercise}[$\star\star$]\label{exo:b2:fragment-orbitals:5}
अमोनिया की प्रथम आयनन ऊर्जा \qty{10.07}{eV} है, मेथेन की
\qty{12.61}{eV}। हर स्थिति में कौन-सा कक्षक इलेक्ट्रॉन खोता है, और अमोनिया को आयनित करना
आसान क्यों है?
\end{exercise}

\begin{exercise}[$\star\star$]\label{exo:b2:fragment-orbitals:6}
रैखिक और मुड़े जल के भरे स्तरों की तुलना कीजिए, और समझाइए कि चार संयोजकता इलेक्ट्रॉनों वाला
\ce{BeH2} जैसा अणु रैखिक क्यों है।
\end{exercise}

\begin{exercise}[$\star\star$]\label{exo:b2:fragment-orbitals:7}
द्वि-स्तरीय परिणाम से कार्बधनायनों
\ce{CH3+}, \ce{CH3CH2+}, \ce{(CH3)2CH+}, \ce{(CH3)3C+} के स्थायित्व का क्रम समझाइए।
\end{exercise}

\begin{exercise}[$\star\star$]\label{exo:b2:fragment-orbitals:8}
आबंध के परितः कोण $\theta$ से ऐंठे दो p कक्षकों का अतिव्यापन
$\cos\theta$ के समानुपाती है। \omterm{def:b2:diatomic-mos:integrals}{अनुनाद समाकल} को अतिव्यापन के समानुपाती मानते हुए,
$\pi$ स्थायीकरण $\theta$ पर कैसे निर्भर करता है? किस कोण पर यह
आधा रह जाता है?
\end{exercise}

\begin{exercise}[$\star\star$]\label{exo:b2:fragment-orbitals:9}
बोरेन \ce{BH3} समतलीय है, अमोनिया पिरामिडी। केंद्रीय परमाणु और तीन हाइड्रोजनों के
\omterm{def:b2:fragment-orbitals:fragment}{अणुखंड कक्षकों} से, हर एक के इलेक्ट्रॉनों के आधार पर, अंतर समझाइए।
\end{exercise}

\begin{exercise}[$\star\star\star$]\label{exo:b2:fragment-orbitals:10}
\ce{H3+} आयन एक समबाहु त्रिभुज है। 1s कक्षकों के लिए $\alpha$ और $\beta$ के साथ,
अतिव्यापन नगण्य मानकर, उसकी तीनों कक्षक ऊर्जाएँ ज्ञात कीजिए (सममित
संयोजन की ऊर्जा $\alpha + 2\beta$ है) और समझाइए कि उसे बाँधने के लिए दो इलेक्ट्रॉन
क्यों पर्याप्त हैं।
\end{exercise}

\begin{exercise}[$\star\star\star$]\label{exo:b2:fragment-orbitals:11}
ऐलिल निकाय \ce{CH2=CH-CH2} के तीन समांतर p कक्षक हैं। नोडों की संख्या से उसके तीनों
$\pi$ कक्षकों की आकृतियों (हर कार्बन पर चिह्न) का अनुमान लगाइए, और बताइए कि
ऐलिल ऋणायन का उच्चतम भरा स्तर कौन-सा है।
\end{exercise}

\begin{exercise}[$\star\star\star$]\label{exo:b2:fragment-orbitals:12}
O--C--O अक्ष के लंबवत 2p कक्षकों से कार्बन डाइऑक्साइड का $\pi$ निकाय
बनाइए: कितने $\pi$ कक्षक, कौन-से आबंधी, अनाबंधी,
प्रतिआबंधी हैं, और उनमें कितने इलेक्ट्रॉन हैं?
\end{exercise}

\section{समस्या: जल की आकृति}

\begin{problem}[{सप्ताहांत समस्या --- ऑक्सीजन और \ce{H2} \omterm{def:b2:fragment-orbitals:fragment}{अणुखंड}, रैखिक
H--O--H के कक्षक, मोड़ने पर खिसकने वाले स्तर, और कूपमान्स सन्निकटन के साथ
प्रकाशइलेक्ट्रॉन स्पेक्ट्रम}]
\label{pb:b2:fragment-orbitals:1}
आँकड़े: जल की ज्यामिति, $r(\ce{O-H}) = \qty{95.8}{pm}$, कोण
$104.48^\circ$; प्रथम आयनन ऊर्जाएँ: \ce{H2O} \qty{12.62}{eV}, \ce{O}
\qty{13.62}{eV}। जल $yz$ तल में है, $z$ समद्विभाजक के अनुदिश।
% ledger: geo:H2O.rOH, geo:H2O.aHOH, ie:H2O, ie:O

\medskip
\noindent\textbf{भाग I --- \omterm{def:b2:fragment-orbitals:fragment}{अणुखंड}।}
\begin{enumerate}
  \item जल में कितने संयोजकता इलेक्ट्रॉन हैं, और किन परमाणुओं से?
  \item हाइड्रोजन 1s कक्षकों के दोनों \omterm{def:b2:fragment-orbitals:symmetry-adapted}{सममिति-अनुकूलित संयोजन} लिखिए।
  \item ऑक्सीजन के कौन-से कक्षक दोनों दर्पण तलों के सापेक्ष सममित हैं?
  \item ऑक्सीजन का कौन-सा कक्षक केवल $xz$ तल के सापेक्ष प्रतिसममित है?
  \item ऑक्सीजन के किस कक्षक का हाइड्रोजन संयोजनों में कोई साथी नहीं है?
  \item कितने \omterm{def:b2:diatomic-mos:lcao}{आण्विक कक्षक} बनते हैं, और कितने भरे हैं?
\end{enumerate}

\medskip
\noindent\textbf{भाग II --- रैखिक H--O--H।}
\begin{enumerate}[resume]
  \item $z$ के अनुदिश रैखिक अणु में कौन-सा संयोजन
        ऑक्सीजन $2\mathrm p_z$ से अन्योन्यक्रिया करता है?
  \item कौन-सा ऑक्सीजन 2s से अन्योन्यक्रिया करता है?
  \item ऑक्सीजन के कौन-से कक्षक अनाबंधी रहते हैं, और किस अपभ्रष्टता के साथ?
  \item आठ संयोजकता इलेक्ट्रॉनों से स्तरों को भरिए।
  \item क्या रैखिक जल \omterm{def:b2:diatomic-mos:magnetism}{अनुचुंबकीय} होता?
  \item उसका उच्चतम भरा स्तर, ऑक्सीजन 2p कक्षक की तुलना में, कहाँ
        होता?
\end{enumerate}

\medskip
\noindent\textbf{भाग III --- मोड़ना।}
\begin{enumerate}[resume]
  \item $z$ के अनुदिश पड़े रैखिक अणु को $yz$ तल में मोड़ा जाता है। कौन-सा
        \omterm{def:b2:diatomic-mos:bonding}{अनाबंधी कक्षक} $h_1 + h_2$ से अतिव्यापित होना आरंभ करता है?
  \item उसके द्वारा वहन किए गए भरे स्तर का क्या होता है?
  \item कौन-सा आबंधी स्तर थोड़ा ऊपर उठता है, और क्यों?
  \item निष्कर्ष निकालिए: जल मुड़ा हुआ क्यों है?
  \item मापे गए कोण की तुलना $90^\circ$ (शुद्ध p आबंधन) और $109.5^\circ$ से कीजिए।
  \item कौन-सा \omterm{def:b2:diatomic-mos:bonding}{अनाबंधी कक्षक} मोड़ने से अप्रभावित रहता है?
\end{enumerate}

\medskip
\noindent\textbf{भाग IV --- प्रकाशइलेक्ट्रॉन स्पेक्ट्रम।}
\begin{enumerate}[resume]
  \item स्पेक्ट्रम कितने संयोजकता बैंड दिखाता है?
  \item कूपमान्स सन्निकटन से, पहला बैंड किस कक्षक से आता है?
  \item वह बैंड संकरा क्यों है, जबकि \omterm{def:b2:diatomic-mos:bonding}{आबंधी कक्षकों} के बैंड चौड़े होते हैं?
  \item जल की प्रथम आयनन ऊर्जा की तुलना ऑक्सीजन
        परमाणु की प्रथम आयनन ऊर्जा से कीजिए, और अंतर समझाइए।
  \item लुइस संरचना के ``दो तुल्य एकाकी युग्मों'' का क्या होता है?
  \item जल के संयोजकता आयनन बैंडों की संख्या और सबसे निचले बैंड की ऊर्जा
        बताइए।
\end{enumerate}
\end{problem}
